\chapter{Builder}
\label{ch:builder}

\section{What will you achieve?}

Builder lets you design an application for your own work when the ready-made apps have no record that fits it, such as an internal request, an equipment register or visit tracking. You first decide what one record represents, then add its fields, arrange its page and lists, and publish it for the team. Later you can add a request form, an approval workflow, documents generated from its data, and a dashboard that tracks it.

This work belongs to the system administrator or anyone with Builder permissions, and it needs a full seat: Builder is not shown to members with an employee seat. Employees use the application after it is published, without needing Builder.

\begin{rcnote}[Example: Horizon]
Horizon wants to receive branch visit requests in one place instead of scattered messages. The system administrator therefore builds a record type called “Visit requests”, where Omar Haddad submits the Aleppo Branch visit request and Sarah Al-Khatib reviews it.
\end{rcnote}

\begin{figure}[H]
\centering
\begin{tikzpicture}[node distance=0.5cm]
\node[stepBox] (a) {Decide what a record is};
\node[stepBox,below=of a] (b) {Add the fields};
\node[stepBox,below=of b] (c) {Arrange page and lists};
\node[stepBox,below=of c] (d) {Rules and access};
\node[stepBox,below=of d] (e) {Review and publish};
\node[stepBox,below=of e] (f) {Show it in navigation};
\node[stepBox,below=of f] (g) {Forms, documents, reports};
\draw[flowsTo] (a) -- (b);
\draw[flowsTo] (b) -- (c);
\draw[flowsTo] (c) -- (d);
\draw[flowsTo] (d) -- (e);
\draw[flowsTo] (e) -- (f);
\draw[flowsTo] (f) -- (g);
\end{tikzpicture}
\rclegend
\caption{Building a new application}
\label{fig:builder-overview-flow}
\shotnote{The figure shows the order this chapter recommends: build the record type first, because everything after it depends on its fields. Then publish it and place it in navigation so the team can reach it, and once it works add forms, documents and reports. The team sees no change before you publish, so try each stage with a real record before moving to the next.}
\end{figure}

\section{How do the Builder tools fit together?}
\label{sec:builder-map}

Builder can look like many tools, but they revolve around one idea: the record type. It is where the data is stored, and every other tool works with it in a different way: a form collects data into it, actions run work on it, templates produce documents from it, and dashboards show its numbers.

\begin{figure}[H]
\centering
\begin{tikzpicture}[node distance=0.5cm]
\node[entity] (type) {Record type};
\node[entityOpt,below right=0.8cm and 0.6cm of type] (forms) {Forms\\collect data};
\node[entityOpt,below=0.8cm of type] (sop) {Workflows\\run the work};
\node[entityOpt,below left=0.8cm and 0.6cm of type] (out) {Templates and dashboards\\produce results};
\draw[dependsOnOpt] (type.south) -- ++(0,-0.35) -| (forms.north);
\draw[dependsOnOpt] (type.south) -- (sop.north);
\draw[dependsOnOpt] (type.south) -- ++(0,-0.35) -| (out.north);
\end{tikzpicture}
\rclegend
\caption{The big picture: everything revolves around the record type}
\label{fig:builder-big-picture}
\shotnote{The figure puts the record type on top because the other tools need it, and all of them are optional. Forms bring data in, workflows run work on it, and templates and dashboards produce documents and numbers from it. A record type the team fills in directly may be enough, and you add a form, template or dashboard when the need appears. So do not start with automation or reports before the record type works with real records.}
\end{figure}

\begin{table}[H]
\centering
\caption{Words you will meet in this chapter}
\label{tab:builder-terms}
\begin{tabularx}{\linewidth}{R{3.6cm} Y}
\rcHeaderRow
\thh{Word} & \thh{What it means, simply}\\
Record type (entity) & A template for a group of similar records, such as “Visit requests”. Each request in it is a record.\\
Field & One piece of information in a record, such as the visit date.\\
Builder workspace & Where you edit a record type; its sections are Fields, Record page, Lists, Rules, Access and more.\\
Draft and publish & Every change is saved as a draft nobody sees, then reaches the team when you publish a new version.\\
Role and verb & A role is a set of permissions given to people; a verb is an allowed action, such as adding or changing.\\
Form & A page of questions an employee or applicant fills in, whose answers may be written into a record type.\\
Internal SOP & A workflow of steps, such as requesting approval or creating a task, that starts on an event you choose.\\
Template & A Word or Excel file with tags that are replaced by record data to create a document.\\
\end{tabularx}
\end{table}

\section{The Builder guide}
\label{sec:builder-pages}

Builder is broad, and each of its tools has many details, so it has a guide of its own whose chapters start from 1. Each chapter explains one tool from start to finish and splits it into short pages, each with a Horizon example, numbered screenshots and steps to check the result. Start with the “Entities” chapter, because the other tools work on the record types you build there, then move on to the other chapters as your work needs.

\begin{rcnote}[Example: Horizon]
Horizon's system administrator starts with the “Entities” chapter to build and publish “Visit requests”, then moves to the “Apps \& Navigation” chapter to place it in the “Branch services” app. When Omar Haddad needs a form to submit the request from, the administrator opens the “Forms” chapter.
\end{rcnote}

\begin{table}[H]
\centering
\caption{Chapters of the Builder guide}
\label{tab:guide-builder-chapters}
\rcZebra
\begin{tabularx}{\linewidth}{R{4.2cm} Y}
\rcHeaderRow
\thh{Chapter} & \thh{What it explains}\\
Entities & Creating a record type and building it in the Builder workspace: fields, record page, lists, rules and access, then reviewing and publishing it.\\
Apps \& Navigation & Arranging the apps and sections members see, and what each member sees first.\\
Forms & Building forms that employees or applicants fill in, linking their answers to record types, then publishing them.\\
Services & Offering forms as services that the public or employees request, and following up requests and their performance.\\
Internal SOPs & Building workflows of steps that start on an event you choose, such as requesting an approval or creating a task.\\
SOPs and Bots & Automations that connect KayanOS with other services and accounts, and bots that work in chat, email and calendar.\\
Documents \& printing & Word, Excel and print templates, and the identities and letterheads that appear on documents.\\
Dashboards \& reports & Dashboards of figures, standalone pages, and AI analysis and reports.\\
\end{tabularx}
\end{table}

“Status templates” and “Task settings” also appear in Builder's “Core” section; the Tasks chapter explains them: the first on page \ref{sec:tasks-feature-statuses}, and the second on page \ref{sec:tasks-feature-settings}.
